
\section*{4. Calcul du \(K^{\bullet}\) d'un fibré de drapeaux}

4.1.
Soient \(S\) un schéma, et \(\underline{E}\) un \(\underline{O}_S\)-Module localement libre de rang \(n\). Soit \(\pi=(p_1,\ldots,p_k)\) une suite d'entiers \(\geq1\), telle que~:
\[
  \sum_{i=1}^k p_i=n.
\]
Rappelons qu'on appelle foncteur des drapeaux de type \(\pi\) de \(\underline{E}\) le foncteur sur la catégorie des \(S\)-schémas, à valeurs dans la catégorie des ensembles, qui à un \(S\)-schéma \(X\xrightarrow{f} S\) associe l'ensemble des suites de quotients de \(f^*(\underline{E})\)~:
\[
  (\underline{L}_{k-1},\ldots,\underline{L}_1),
\]
où, pour tout \(i\), \(\underline{L}_i\) est un quotient de \(\underline{L}_{i+1}\), localement libre de rang \(p_1+\ldots+p_i\).

Un tel foncteur est représentable par un \(S\)-schéma \(\mathbf D_\pi(\underline{E})\), qui est muni d'une suite de quotients canonique de l'image inverse de \(\underline{E}\), localement libres de rang \(p_1,p_1+p_2,\ldots\)

Si \(\pi=(1,n-1)\), alors \(\mathbf D_\pi(\underline{E})=\mathbf P(\underline{E})\).

Si \(\pi=(r,n-r)\), alors \(\mathbf D_\pi(\underline{E})=\mathbf G_{r,n}(\underline{E})\), grassmannienne de type \((r,n-r)\) de \(\underline{E}\).

Si tous les \(p_i\) sont égaux à \(1\), on trouve le fibré en drapeau usuel de \(\underline{E}\), noté \(\mathbf D(\underline{E})\).

\textbf{Lemme 4.2.} \uline{Soient} \(S\) \uline{un schéma}, \(\underline{E}\) \uline{un} \(\underline{O}_S\)\uline{-Module localement libre de rang} \(n\), \(\pi=(p_1,\ldots,p_i,\ldots,p_k)\), \uline{avec} \(p_1+\ldots+p_k=n\) ; \uline{soient} \(p',p''\) \uline{deux entiers} \(\geq1\), \uline{tels que} \(p'+p''=p_i\) ; \uline{posons} \(\pi'=(p_1,\ldots,p_{i-1},p',p'',p_{i+1},\ldots,p_k)\), \(\pi''=(p',p'')\). \uline{Soient} \(\underline{L}_1,\ldots,\underline{L}_{k-1}\) \uline{les quotients canoniques de l'image inverse de} \(\underline{E}\) \uline{sur} \(\mathbf D_\pi(\underline{E})\), \uline{et soit} \(\underline{F}\) \uline{le noyau de} \(\underline{L}_i\to \underline{L}_{i-1}\). \uline{Alors on a un} \(S\)-\uline{isomorphisme canonique}~:
